% Response to reviewers for the revised PoolParty manuscript.
\documentclass[11pt]{article}

\usepackage[margin=1in]{geometry}
\usepackage[T1]{fontenc}
\usepackage[utf8]{inputenc}
\usepackage{lmodern}
\usepackage{xcolor}
\usepackage{enumitem}
\usepackage{parskip}
% Match main.tex's sn-vancouver-num class option: bracketed, numerical
% citations that are sorted and compressed by natbib.
\usepackage[numbers,sort&compress]{natbib}
\usepackage{xr-hyper}
\usepackage[hidelinks]{hyperref}
% Labels from both documents are imported unprefixed, so quoted blocks can be
% pasted verbatim from main_bmc.tex and si.tex. Safe because response.tex defines
% no labels of its own and the two label sets are disjoint. The "nocite"
% argument suppresses import of \bibcite entries, which would otherwise clash
% with this document's own citation numbering.
\externaldocument[][nocite]{main}
\externaldocument[][nocite]{si}

\definecolor{responseblue}{HTML}{0070C0}
\definecolor{noticegray}{HTML}{555555}
% Must match \quoted in main_bmc.tex and si.tex.
\definecolor{quotedtext}{HTML}{A00020}
\newenvironment{response}{\par\color{responseblue}}{\par}
\newenvironment{msquote}
  {\begin{quote}\color{quotedtext}\noindent\ignorespaces}
  {\end{quote}}
\providecommand{\bibcommenthead}{}
% Flat table of contents: every entry is one flush-left, dotted line.
\makeatletter
\renewcommand*{\l@section}{\@dottedtocline{1}{0pt}{0pt}}
\makeatother
\newcommand{\tocline}[1]{\addcontentsline{toc}{section}{#1}}
\setlist[enumerate,1]{leftmargin=*,itemsep=1.2\baselineskip,topsep=\baselineskip}
\newcommand{\reviewersection}[1]{%
  \phantomsection
  \tocline{#1}%
  \bigskip\par\noindent{\large\bfseries #1}\par\medskip\hrule\medskip}

\begin{document}

\begin{center}
{\Large\bfseries Response to Reviewers}\\[0.6em]
{\large PoolParty: streamlined design of DNA sequence libraries in Python}\\[0.8em]
Zhihan Liu, Aidan Cordero, Justin B. Kinney\\[0.7em]
\end{center}

\renewcommand{\contentsname}{Contents}
\tableofcontents
%\clearpage

\reviewersection{Overview}

\begin{response}
We thank the Editor and Reviewers for their thoughtful critiques. The manuscript has been improved substantially as a result of this feedback. Additional edits have also been made throughout to improve clarity. We provide a point-by-point response to each critique below.
\end{response}

\reviewersection{Editor Comments}

In general, the reviewers had positive impressions of PoolParty, and it can
likely be revised to a state that will be acceptable. Please carefully address
each point raised by the reviewers. It's especially important to provide
benchmarking data, comparisons to existing tools, and not to overstate the
benefits the tool may provide or fail to discuss limitations. Reviewer 1
recommends comparing PoolParty to 2--3 other comparable tools. However, rather
than arbitrarily picking two or three other tools for comparisons, you should
use a well-defined process for picking which tools, and how many, to include in
comparisons. Perhaps 2--3 is adequate; however, that may be too few tools for
comparisons. The comparisons should include, at a minimum, benchmarking (e.g.,
runtime, memory usage), program features, and similarities/differences in
program outputs.

\begin{response}

We have expanded the Background section to discuss in more depth the limitations of competing tools that PoolParty overcomes:
  \begin{msquote}
    Existing tools can help with parts of this process. General-purpose sequence toolkits \citep{Cock2009df, Pereira2015wj, Klie2023kg} can manipulate sequences but have no notion of a variant library as a structured object. Consequently, users must write custom scripts to handle the combinatorial logic themselves. Several other tools automate library design or sequence perturbation for particular applications. For example, VaLiAnT \citep{Barbon2022th} generates mutational scanning libraries of both coding and noncoding regions; supported mutation types include single-nucleotide substitutions, amino acid substitutions, and in-frame deletions. For MPRAs, MPRAnator \citep{GeorgakopoulosSoares2017gb} provides modules for motif placement, SNP design, and negative control generation \citep[see also][]{Ghazi2018aa}. For in silico sequence experiments, tangermeme \citep{Schreiber2025nd} supports sequence perturbation and analysis of genomic deep learning models. However, each of these tools is targeted to specific applications and, as a result, no individual tool can generate the full range of library designs in common use. For example, VaLiAnT does not support regulatory sequence libraries containing random combinations of transcription factor binding sites, while MPRAnator and tangermeme do not support codon-aware mutagenesis of protein-coding sequences.
  \end{msquote}

The revised Discussion now better qualifies how PoolParty compares to existing tools:
  \begin{msquote}
    PoolParty complements existing library-design tools such as VaLiAnT \citep{Barbon2022th}, MPRAnator \citep{GeorgakopoulosSoares2017gb}, and tangermeme \citep{Schreiber2025nd} by providing a single framework for DMS libraries, MPRA libraries, and library designs that combine elements of both. \ldots{} Additional comparisons between PoolParty and existing tools are provided in Supplemental Information, Figs.~\ref{fig:valiant-comparison} and \ref{fig:tangermeme-comparison} and Tables~\ref{tab:tools} and \ref{tab:capabilities}.
  \end{msquote}
The revised Discussion now also describes the limitations of PoolParty explicitly:
  \begin{msquote}
    PoolParty was designed to solve a specific computational problem: streamlining the design of DNA sequence libraries for a wide range of experimental applications and in silico studies. As such, it has several limitations. First, PoolParty operates on sequences, not genomic loci or transcript models. It can accept genomic coordinates as input to extract sequences from FASTA files, but does not interpret transcript annotations (e.g., in GFF/GTF files) or keep track of genomic coordinates through subsequent design steps. This may make PoolParty less helpful for certain genome-editing applications (e.g., as compared to VaLiAnT). PoolParty is also not a sequence optimization tool; it does not design sequences that have optimized codon usage \citep{Swainston2017rb}, that are subject to complex synthesis constraints \citep{Zulkower2020jk}, or that have optimal activity as assessed by machine learning models \citep{Schreiber2025ledidi}. PoolParty also does not address the problem of physically constructing DNA sequence libraries, e.g., by designing mutagenic primers \citep{Hiraga2021yg}. We expect, however, that PoolParty can in many cases be productively paired with complementary tools that do address these challenges.
  \end{msquote}

We have also added a Supplemental Information document that compares PoolParty to competing tools (Tables \ref{tab:tools} and \ref{tab:capabilities}, Figures \ref{fig:valiant-comparison} and \ref{fig:tangermeme-comparison}). The approach used to select these tools is described in the Supplemental Methods section of that document:
  \begin{msquote}
    A literature search identified 13 tools with substantial functional overlap with PoolParty; these are listed in Table \ref{tab:tools}. Table \ref{tab:capabilities} compares the specific features and capabilities of the 7 tools most similar in purpose to PoolParty. Of these 7, we judged VaLiAnT \citep{Barbon2022th}, MPRAnator \citep{GeorgakopoulosSoares2017gb}, and tangermeme \citep{Schreiber2025nd} to be the closest competitors. MPRAnator is a web application that, as of this writing, is no longer accessible. Figs.~\ref{fig:valiant-comparison} and \ref{fig:tangermeme-comparison} provide explicit comparisons of the code/inputs required to construct identical libraries using PoolParty versus VaLiAnT (Fig.~\ref{fig:valiant-comparison}) or PoolParty versus tangermeme (Fig.~\ref{fig:tangermeme-comparison}). We note that neither VaLiAnT nor tangermeme is capable of generating all three libraries used in the benchmarking study of PoolParty.
  \end{msquote}

  The Supplemental Information also reports benchmarking experiments that quantify the compute and memory requirements of PoolParty on the library designs featured in the main text (Fig.~\ref{fig:benchmark-scaling} and Table \ref{tab:benchmark-summary}). We did not extend these benchmarking experiments to other tools for two reasons. First, none of the other tools was capable of generating libraries with the same three designs. Second, our paper does not claim that PoolParty has any execution speed or memory usage advantages over other methods. Rather, the competitive advantage of PoolParty is its flexibility, ease of use, and the per-sequence metadata it provides. 
\end{response}

\reviewersection{Reviewer 1}

\begin{enumerate}
\item In the GB1 example, the library exceeds 540,000 sequences, yet the paper
provides no runtime or memory consumption benchmarks. Readers cannot assess the
tool's practical usability across different library scales. The authors should
supplement benchmarking data, including runtime and peak memory usage across
libraries of increasing sizes (e.g., 1K, 10K, 100K, 500K sequences), and specify
the test environment configuration.

\begin{response}
This is a good suggestion. As described above, the new Supplemental Information document (Fig.~\ref{fig:benchmark-scaling} and Table \ref{tab:benchmark-summary}) reports benchmarking tests of runtime and peak memory usage using libraries of increasing size that match the designs used for Figs.~\ref{fig:figure2}, \ref{fig:figure3}, and \ref{fig:figure4} of the main text. The test environment is reported in the Supplemental Methods section of that document. 
  \begin{msquote}
    Figure~\ref{fig:benchmark-scaling} and Table \ref{tab:benchmark-summary} present a benchmarking study of PoolParty using the three library designs featured in Figs.~2, 3, and 4 of the main text. Benchmarking experiments were run under Ubuntu 22.04.3 LTS in Windows Subsystem for Linux 2 on an Intel Core Ultra 9 185H workstation with 31~GB RAM allocated to WSL2, using CPython 3.12.2, PoolParty 0.1.0, and StateTracker 0.1.0. Timings are single-threaded wall-clock measurements of \texttt{generate\_library}, reported as mean $\pm$ SD over five replicate runs. Peak memory is the maximum resident set size reported by \texttt{getrusage}, with each data point measured in a fresh process.
  \end{msquote}
\end{response}

\item The paper highlights design cards as a core feature, emphasizing their use
as covariates in downstream analyses. However, it does not evaluate whether
design card variables provide significant incremental predictive value beyond
sequence-derived features. The authors should either supplement a comparative
experiment quantifying the added information gain of design cards, or explicitly
state in the Discussion that this validation is planned as future work.

\begin{response}
The purpose of these design cards is simply to provide users with easy access to the design choices used to construct each sequence. We do not mean to claim that this information can increase the predictive performance of sequence-to-function models. We have clarified the intended role of design cards in the Discussion:
  \begin{msquote}
    Design cards are thus a bookkeeping device that allows users to avoid reverse-engineering sequence architecture from raw sequence when investigating how this architecture affects the output of a sequence-to-function model.
  \end{msquote}
\end{response}

\item The introduction mentions VaLiAnT and MPRAnator, noting their limitations,
but does not provide a systematic comparison between PoolParty and these tools.
The authors should provide a table comparing PoolParty with 2--3 existing tools
on core features, and demonstrate code conciseness in at least one common use
case, to help readers objectively assess PoolParty's advantages.

\begin{response}
This is a good suggestion. Tables \ref{tab:tools} and \ref{tab:capabilities} in the new Supplemental Information document provide systematic comparisons between PoolParty and other competing tools; Figs.~\ref{fig:valiant-comparison} and \ref{fig:tangermeme-comparison} further provide explicit comparisons of the code/inputs required to construct identical libraries using PoolParty versus VaLiAnT (Fig.~\ref{fig:valiant-comparison}) or PoolParty versus tangermeme (Fig.~\ref{fig:tangermeme-comparison}). Supplemental Methods describes the specific process we used for choosing the tools featured in these comparisons. Please note that MPRAnator (which runs through a web interface) is no longer functional, and so we were unable to carry out a similar comparison between it and PoolParty.
\end{response}

\item Provide benchmarking data on runtime and memory usage across libraries of
varying sizes.

\begin{response}
This is addressed in the response to Comment 1 above. 
\end{response}

\item Include a table comparing features with 2--3 existing tools.

\begin{response}
This is addressed in the response to Comment 3 above. 
\end{response}

\item To further contextualize this research within the existing academic
landscape and strengthen its theoretical support, it is recommended that the
authors supplement several relevant representative references.

[1] \url{https://doi.org/10.1093/nsr/nwaf495}

[2] \url{https://doi.org/10.1038/s41467-026-72882-y}

[3] \url{https://doi.org/10.1186/s13059-025-03606-6}

[4] \url{https://doi.org/10.1016/j.patcog.2025.112078}

\begin{response}
It is unclear to us how these papers are relevant to PoolParty or the more general problem of sequence library design. 
\begin{itemize}
  \item The first paper, ``Multimodal pre-training models of molecular representation for drug discovery'' by \citet{Cui2025sg}, is a review article that focuses on applications of AI in drug discovery, with a specific focus on the design of small-molecule drugs. It does not review algorithms or software available for designing sequence libraries. 
  \item The second paper, ``CAPTAIN: a multimodal foundation model pretrained on co-assayed single-cell RNA and protein'' by \citet{Ji2026ct}, describes a foundation model that aims to link single-cell transcriptomic measurements to protein expression profiles. It does not introduce any strategies for sequence library design.
  \item The third paper, ``A survey of sequence-to-graph mapping algorithms in the pangenome era'' by \citet{Wang2025mp}, is a review article discussing the computational methods that have been developed to map sequences to pangenome graphs. It does not review algorithms or software available for designing sequence libraries. 
  \item The fourth paper, ``Multi-hop graph structural modeling for cancer-related circRNA-miRNA interaction prediction'' by \citet{Wei2026mh}, describes a computational method (called GraCMI) for predicting interactions between circRNAs and miRNAs. It does not address the problem of sequence library design.
\end{itemize}
We have therefore decided not to cite these papers in the manuscript. 
\end{response}

\item Either provide empirical analysis of design cards' incremental
contribution to surrogate modeling performance, or explicitly discuss this as
future work.

\begin{response}
This is addressed in the response to Comment 2 above. 
\end{response}
\end{enumerate}

\reviewersection{Reviewer 2}

\begin{enumerate}
\item Additional statistical readouts describing pool generation would be
useful. For example, how many unique sequences vs duplicates were produced in
each pool (i.e. out of the universe of all permutations that meet the
specifications)? How unique are the sequences (min/max/average hamming distance)?
How frequent are homopolymers? Etc.

\begin{response}
This is a good suggestion. To make this information readily available to users, we have added a \texttt{stats} method to the Pool class. This method computes summary statistics including the size of the design space (when this is known), the number of unique/duplicated sequences, the pairwise Hamming distances between sequences, the frequencies of homopolymer occurrences, and so on. We have also added a short paragraph to the Sequence metadata subsection describing this method.
  \begin{msquote}
    PoolParty also provides a \texttt{stats} method for summarizing a DNA Pool, including its unique and duplicate sequence counts, pairwise Hamming distances, and homopolymer frequencies.
  \end{msquote}
\end{response}

\item The authors disclose that LLMs were used to develop, debug, and document the code. Did the authors also write a test harness, and what edge cases were nominated by the LLMs for testing vs what edge cases were nominated by the human developers for testing?

\begin{response}
The PoolParty codebase includes over 3{,}000 tests that can be run using pytest. These include a combination of LLM-written tests and tests derived from human-written use cases. In our experience, the LLM coding agents were quite good at testing mechanical aspects of the code (e.g., invalid inputs and boundary conditions). The human-written use cases, however, were essential for testing the types of inputs and requests likely to be encountered in real biological applications. We did not track the origin of each test, however, and so cannot precisely quantify how many came from each source.  
\end{response}

\item The backward/forward pass steps are not well described in the manuscript,
and figure 1b is hard to interpret. I suggest adding a supplemental figure
showing how these passes are performed. For example, the text explains that the
backward pass happens as information is passed to the root node -- does this
mean that the state is passed from `Mutagenize' to `Pool A' in or first from
`Pool Final' to `Stack'? And if so, what information is passed? Perhaps showing
each step in an expanded figure or labeling arrows with the order would be more
useful.

\begin{response}
This is a good suggestion. Our new Supplemental Information document includes a figure, Fig.~\ref{fig:sequence-generation}, that breaks down the forward and backward pass computations in much more detail. In this particular case, the backward pass involves `Pool Final' (the root node) first passing its state (an integer encoding which specific sequence to generate) to `Stack', which then decomposes its state and passes the results to `Pool B' and `Pool C', and so on. 
\end{response}

\item A comparison to the pools produced by other tools such as VaLiAnT or
MPRAnator would provide additional confidence in the ability of PoolParty to
produce sequence pools as expected. Metrics such as the number of pool elements that overlap and an investigation of any unique pool elements would ensure that the generate pools are accurate.

\begin{response}
Figs.~\ref{fig:valiant-comparison} and \ref{fig:tangermeme-comparison} in Supplemental Information now compare the code/inputs needed to create identical libraries with PoolParty and VaLiAnT (Fig.~\ref{fig:valiant-comparison}) or tangermeme (Fig.~\ref{fig:tangermeme-comparison}). We attempted to make a similar figure for MPRAnator, but that tool uses a web interface that is currently inoperative. 
\end{response}

\item Additional operations to improve library content would be useful. E.g. to maximize hamming distance, remove sequences containing restriction enzyme sites, remove homopolymers, etc. would make the tool more useful in generating ready-to-use pools.

\begin{response}
The original manuscript failed to mention that PoolParty provides a \texttt{filter} Operation that allows users to remove sequences that do not satisfy user-specified criteria. This Operation includes a number of built-in criteria, including GC content, homopolymer content, sequence complexity, and restriction enzyme recognition sites. To clarify this capability, we have added the following sentence to the Pools and Operations subsection of the main text:

\begin{msquote}
Sequences can be filtered based on sequence length, GC content, homopolymer length, sequence complexity, restriction enzyme recognition sites, or other user-defined criteria.
\end{msquote}
\end{response}

\begin{response}
We have also clarified PoolParty's limitations in the Discussion section, in particular the fact that it is not designed to generate libraries that are optimized for specific metrics.

\begin{msquote}
PoolParty is also not a sequence optimization tool; it does not design sequences that have optimized codon usage \citep{Swainston2017rb}, that are subject to complex synthesis constraints \citep{Zulkower2020jk}, or that have optimal activity as assessed by machine learning models \citep{Schreiber2025ledidi}.
\end{msquote}
\end{response}

\item The name PoolParty may be confused with the existing PoolParty analysis
softwre (\url{https://doi.org/10.1111/1755-0998.12784}).

\begin{response}
This is a fair concern, but we have decided to retain the name PoolParty. Our package already uses the name \texttt{poolparty} on PyPI and ReadTheDocs, and we have already announced this software in a preprint. Moreover, the previously published PoolParty is a pipeline for aligning and analyzing pooled-sequencing data, a use case that we believe is different enough from library sequence generation to avoid confusion with our package. 
\end{response}

\end{enumerate}

\reviewersection{Reviewer 3}

\begin{enumerate}
\item I think that the emphasis made on PoolParty being a universal tool that
can ``replace'' (line 256) existing code/tools needs to be qualified more. The
benefits of PoolParty are that it is straightforward to use, interpretable,
modular, with a nice `design card' tracking and graphed interpretation -- I
think this is particularly useful for ML and in silico work.

However, whist it does streamline design processes, the package operates on
input sequences rather than genomic loci co-ordinates/transcript so in terms of
wet-lab experiments it is most useful for fundamental biology cDNA DMS and MPRA
assays that are episomal (i.e. from plasmid / using reporters) rather than MAVEs
that rely upon genome editing (which is becoming the gold standard for clinical
variant work). It does not intuitively and natively support reference genome
coordinates, transcript models, exon/intron structures, alternative isoforms,
or genome assembly awareness (as far as I can tell); VaLiAnT does do this. Base
sequences into which variants are installed are represented/input as sequence
string (pasted by user) rather than standardized co-ordinate or HGVS notation
(this limits interpretation/ingress of variants from clinical and genomics
resources). If PoolParty is presented as intuitive and adaptable rather than a
universal replacement to alternative tools and code this is less of a concern
and may allow researchers to more readily select which software to invest time
into learning/using.

\begin{response}
This is a fair critique: PoolParty does not currently operate on annotated genomic loci, nor does it accept HGVS-formatted input. We have removed language suggesting that PoolParty is meant to replace tools that do, and have revised the Discussion section to clarify the intended scope of PoolParty. 
\begin{msquote}
  PoolParty complements existing library-design tools such as VaLiAnT \citep{Barbon2022th}, MPRAnator \citep{GeorgakopoulosSoares2017gb}, and tangermeme \citep{Schreiber2025nd} by providing a single framework for DMS libraries, MPRA libraries, and library designs that combine elements of both. ... Additional comparisons between PoolParty and existing tools are provided in Supplemental Information, Figs.~\ref{fig:valiant-comparison} and \ref{fig:tangermeme-comparison} and Tables~\ref{tab:tools} and \ref{tab:capabilities}. 

  PoolParty was designed to solve a specific computational problem: streamlining the design of DNA sequence libraries for a wide range of experimental applications and in silico studies. As such, it has several limitations. First, PoolParty operates on sequences, not genomic loci or transcript models. It can accept genomic coordinates as input to extract sequences from FASTA files, but does not interpret transcript annotations (e.g., in GFF/GTF files) or keep track of genomic coordinates through subsequent design steps. This may make PoolParty less helpful for certain genome-editing applications (e.g., as compared to VaLiAnT). ...
\end{msquote}
In the new Supplemental Information document, we have also added two tables that explicitly compare the features and capabilities of PoolParty with existing tools (Tables~\ref{tab:tools} and \ref{tab:capabilities}). 
\end{response}

\item Molecular consequence beyond codon change for missense is not reported as
far as I can tell, VaLiAnT does this to some extent (e.g. stop-gained,
synonymous, inframe deletion, etc, PoolParty may do this also, but I couldn't
see this after running a `replacement\_scan'?) but generally researchers use VEP
on their tool outputs to achieve extensive annotation. VEP requires standard
annotation or format (e.g. VCF) for input rather than sequence strings -- if
this is possible for PoolParty (the production of a VCF or a format that VEP
will accept) then one could obtain extensive annotations (e.g. SpliceAI
predictions, EVE, AlphaMissense, PolyPhen, Consequence, genomic coordinates,
Clinvar and gnomAD identifiers etc etc) which I think would greatly increase the
uptake of the tool by researchers in the MAVE community that are more
clinically/genomically minded -- if this is possible I would state this/provide
an easy means to generate VEP input format from PoolParty outputs.

\begin{response}
This is correct: PoolParty does not annotate the molecular consequences of mutations, as it does not return an annotation layer with each sequence. However, if users wish to generate sequences carrying a specific type of mutation and then fold those sequences into a larger library, there are two ways to track which specific sequences received which types of mutations: through the design card provided with each sequence, and through each sequence's name.

We appreciate the suggestion to support VCF format. To this end, we have implemented a new method, \texttt{from\_vcf}, which extracts a specified sequence window and inserts the alternate alleles into that sequence. We do not, however, see a simple way of exporting libraries in VCF format, as the sequences generated by PoolParty are often not closely related to any specific wild-type sequence. The best option in this case would be for users to take the sequences generated by PoolParty, align them to a reference, and use the results to generate a corresponding VCF record. We believe that this last step is best carried out using other available software. 
\end{response}

\item The authors note limitations at the end of the discussion. I think
updating this section to address that PoolParty is mostly of benefit for DMS and
MPRA rather than genome editing MAVE assays (e.g. not as useful for Saturation
Genome Editing HDR libraries and Saturation Prime Editing pegRNA libraries)
which require extensive and specific annotation/standardized nomenclature and
consideration of transcript (generally through ingress of GTF/GFF files) and
exon/intron boundaries (including split codons between exons, again through
GTF/GFF boundary annotations).

Alternatively/additionally the authors might wish to highlight that the tool is
agnostic to species/genome/transcript and present it as a more intuitive
sequence generation tool which is particularly useful for ML
testing/benchmarking which is well tracked with the design cards and
interpretable/iteratively mutable through graph interrogation. This builds on
the comment that generally annotations are variant specific (line 56-58), with
the emphasis being that PoolParty serves as a useful tool to manifest oligo
libraries as structured objects (less of a concern for wet-lab investigations
in my opinion, but much more salient for ML based studies on raw sequence out of
genomic context).

\begin{response}
We have revised the Discussion section to clarify these issues:

\begin{msquote} 
PoolParty operates on sequences, not genomic loci or transcript models. It can accept genomic coordinates as input to extract sequences from FASTA files, but does not interpret transcript annotations (e.g., in GFF/GTF files) or keep track of genomic coordinates through subsequent design steps. This may make PoolParty less helpful for certain genome-editing applications (e.g., as compared to VaLiAnT).
\end{msquote}
\end{response}

\item Line 12 -- there is not a ``lack'' of purpose-built software, replace
with ``scarcity'' perhaps.

\begin{response}
Fair point. We have replaced ``lack'' with ``scarcity'' in the Abstract. The revised sentence reads:

\begin{msquote}
Designing these libraries, however, remains tedious and error-prone due to the scarcity of purpose-built software.
\end{msquote}
\end{response}

\item Line 51-53 - stated that VaLinAnT is assay specific, then say used for DMS
and saturation genome editing two different assays (SGE is a type of DMS, but
VaLinAnT is used for SGE and cDNA DMS which are quite different assays).

\begin{response}
We have replaced ``assay-specific'' with more specific language; the revised background text reads:

\begin{msquote}
Several other tools automate library design or sequence perturbation for particular applications. For example, VaLiAnT \citep{Barbon2022th} generates mutational scanning libraries of both coding and noncoding regions; supported mutation types include single-nucleotide substitutions, amino acid substitutions, and in-frame deletions.
\end{msquote}
\end{response}

\item Line 66 -- ``expensive computation'' this is only true for truly massive
sequence generation -- most full gene DMS/saturation libraries of $>$20,000
variants take $<$5 minutes on a standard laptop CPU (true of Valiant) -- this
statement would be true for many permutations of sequence for use in ML (e.g.
pairwise, full genome CDS on GPU) -- I would qualify this point on cost.

\begin{response}
Indeed, most libraries designed for wet-lab experiments can be generated in seconds to minutes on a standard CPU, as the reviewer notes. We have therefore made the wording more specific:

\begin{msquote}
Importantly, no sequences are generated until the user requests them, and it is simple to generate small test sets of sequences. This allows users to explore multiple design options without having to generate complete libraries at each iteration.
\end{msquote}
\end{response}

\item The use cases of protein coding, MPRA promoter and splice are good examples
(line 68).

\begin{response}
Thanks!
\end{response}

\item Figure 2 B and Figure 3 B -- is it necessary helpful to put these code
blocks here? I'm not sure how useful this is for readership.

\begin{response}
The purpose of these panels is to support our claim that the PoolParty API is simple to use. We believe that showing the code explicitly, together with the output the code generates, is the best way to support this claim, and so we have opted to retain these panels. 
\end{response}

\item Line 194 -- `most abundant codon' -- this needs to be defined, is this
codon usage (I think not) or instance in the input sequence?

\begin{response}
We have clarified the text:

\begin{msquote}
We note that \texttt{mutagenize\_orf} supports multiple codon mutation types; this example uses \texttt{missense\_only\_first}, which selects the most frequently used codon in the human genome for each target amino acid.
\end{msquote}
\end{response}

\item Line 235 -- SpliceAI is defined here for the first time but is introduced
before this in section 2.8 -- I would move the `a deep learning model of mRNA
splicing' to the earlier section.

\begin{response}
The term SpliceAI is now defined in the earlier Implementation subsection, where this term is first used:

\begin{msquote}
SpliceAI \citep{Jaganathan2019ps} is a genomic deep learning model for splice-site prediction in humans. We used the five-model ensemble to predict the canonical 5'ss probability using $\pm$5~kb of genomic context (GRCh38).
\end{msquote}
\end{response}

\item Line 238 ``varying strength'' -- maxentscan in legend, I would define this
is the main body text too.

\begin{response}
We have clarified that cryptic splice-site strength refers to the score assigned by MaxEntScan. The revised sentence in Results reads:

\begin{msquote}
Using the 5'ss of SMN2 exon~7 as the canonical site, we used PoolParty to construct a library in which cryptic 5'ss 9-mers of varying strength
(MaxEntScan score) were inserted at varying positions on either side of the canonical 5'ss (see Implementation).
\end{msquote}
\end{response}

\item Line 248 -- Fig.4G is introduced before 4E-F, perhaps generally introduce
Fig.4. then this could stay as is.

\begin{response}
We have moved the cross-validation result to the end of the paragraph so that
the Figure 4 panels are introduced in order. 
\end{response}

\item Line 256 -- ``replaces the ad hoc'' is too bold and a little presumptive as
for major points above-- needs clarification/qualification. Also, Line 274 --
``in contrast''-- I think these claims are too strong -- valiant can do DMS and
MPRA as well as SGE/SPE -- I would tone these statements down and/or emphasise
what is unique to PoolParty.

\begin{response}
This is fair criticism; we have revised both statements:

\begin{msquote}
We have presented PoolParty, a Python package for designing oligonucleotide libraries. PoolParty facilitates this task by representing libraries as directed acyclic graphs (DAGs) that users can construct using a simple API.
\end{msquote}

\begin{msquote}
PoolParty complements existing library-design tools such as VaLiAnT \citep{Barbon2022th}, MPRAnator \citep{GeorgakopoulosSoares2017gb}, and tangermeme \citep{Schreiber2025nd} by
providing a single framework for DMS libraries, MPRA libraries, and library designs that combine elements of both.
\end{msquote}
\end{response}

\item Line 271 -- ``in the spliceAI example'' -- the one in section 2.8 or 3.3?
(3.3 I imagine, but best to state).

\begin{response}
We have clarified that this refers to the SpliceAI surrogate modeling example in the Results:

\begin{msquote}
For example, in the SpliceAI analysis of Fig.~\ref{fig:figure4}, the cryptic splice-site strength and position recorded in each design card provided the covariates used for surrogate modeling.
\end{msquote}
\end{response}

\item The code is very well written and documented, and it is noted by the
authors that Anthropic models have been used to do this/help which I think is
sensible, editorially this may be something to confirm is acceptable at
Bioinformatics (as reviewers were requested not to use these tools, but this is
likely due to confidentiality).

\begin{response}
BMC Bioinformatics policy permits LLM use when disclosed. We have also clarified the LLM disclosure statement in Implementation:

\begin{msquote}
The authors used large language models (primarily Claude Opus 4.5 and 4.6, Anthropic) to assist with code development, debugging, and documentation writing. The manuscript was written entirely by the authors, with LLM assistance only at the copy-editing stage. The figures were also prepared by the authors; LLMs were used only to help write the scripts that generate the plots. The authors affirm that they are fully responsible for the content of the manuscript, the codebase, and the documentation. 
\end{msquote}
\end{response}

\item The name `PoolParty' is great.

\begin{response}
Thank you!
\end{response}
\end{enumerate}

\phantomsection
\tocline{References}
\bibliographystyle{sn-vancouver-num}
\bibliography{references}

\end{document}
